
%(BEGIN_QUESTION)
% Copyright 2014, Tony R. Kuphaldt, released under the Creative Commons Attribution License (v 1.0)
% This means you may do almost anything with this work of mine, so long as you give me proper credit

Søk i kapittel 1 (``DeltaV System Overview'') i manualen ``Getting Started with your DeltaV Digital Automation System'' (dokument D800002X122, mars 2006) og svar på følgende spørsmål:

\vskip 10pt

Basert på det enkle nettverksdiagrammet på side 1-1: hvilken type digitalt nettverk kobler PC-arbeidsstasjon(er) til DCS-kontrolleren?

\vskip 10pt

Hva er {\it System Identifier} i et DeltaV-system, og hva brukes den til?

\vskip 10pt

Flere programverktøy brukes ofte i drift og vedlikehold av DeltaV-styresystemer: {\it DeltaV Explorer}, {\it DeltaV Operate Run}, {\it DeltaV Operate Configure}, {\it Control Studio} og {\it DeltaV Books Online}. Beskriv hva hvert verktøy brukes til.



\vskip 20pt \vbox{\hrule \hbox{\strut \vrule{} {\bf Forslag til sokratisk diskusjon} \vrule} \hrule}

\begin{itemize}
\item{} Et kjennetegn ved DCS-systemer er {\it maskinvare-redundans}. Pek på konkrete redundante funksjoner i Emerson DeltaV DCS.
\item{} Gå til en DeltaV-arbeidsstasjon og åpne programverktøyene som er listet, og utforsk funksjonene. {\it Ikke ``download'' eller ``save'' noe som endrer DCS-konfigurasjonen -- bare utforsk og observer!}
\end{itemize}

\underbar{file i00811}
%(END_QUESTION)





%(BEGIN_ANSWER)


%(END_ANSWER)





%(BEGIN_NOTES)

Emerson DeltaV DCS bruker redundant Ethernet-kabling for å koble PC-arbeidsstasjoner til kontrollere. Diagrammet med to ``hubs'' som nettverkskabler kobles til (side 1-1) indikerer Ethernet. Redundante prosessorer støttes også.

\vskip 10pt

System Identifier er en maskinvarenøkkel brukt for autentisering av DeltaV-nettverket (side 1-2). Den kobles til enten PC-ens parallellport eller en USB-port.

\vskip 10pt

{\it DeltaV Explorer} brukes til å utforske og redigere komponentkonfigurasjon i et DeltaV-system (side 1-9). {\it DeltaV Operate Run} er operatørgrensesnittet for drift av systemet (side 1-15). {\it DeltaV Operate Configure} er konfigurasjonsmodusen der operatørbilder bygges og redigeres (side 1-10). {\it Control Studio} brukes til å bygge og konfigurere reguleringsstrategiene i systemet, altså programmeringsgrensesnittet for DeltaV (side 1-8). {\it DeltaV Books Online} er et digitalt bibliotek med referansemanualer, inkludert denne ``Getting Started''-manualen.





\vfil \eject

\noindent
{\bf Forberedende quiz:}

Programvaren {\it Control Studio} som følger med Emerson DeltaV DCS brukes vanligvis til hvilket formål:

\begin{itemize}
\item{} Lage grafiske ``picture''-skjermer for operatører
\vskip 5pt 
\item{} Idriftsette nye kontrollernoder på DeltaV-nettverket
\vskip 5pt 
\item{} La operatører bytte mellom Auto- og Manuell-modus
\vskip 5pt 
\item{} Kalibrere prosess-transmitteren for korrekt måling
\vskip 5pt 
\item{} Kalibrere reguleringsventilen for korrekt slaglengde
\vskip 5pt 
\item{} Konfigurere funksjonsblokker som utgjør en reguleringsmodul
\vskip 5pt 
\item{} Beregne installert kostnad for DeltaV-systemet
\end{itemize}


%INDEX% Reading assignment: Emerson DeltaV "Getting Started" manual (Chapter 1)

%(END_NOTES)
